\labsection{6}{Threads, mutexes, and resource safety}{sec:lab06}
\begin{labproject}{Observe a race condition, protect a shared counter, and trace Banker's safety test}
\textbf{Suggested time:} 110--120 minutes \hfill
\textbf{Level:} intermediate\par
\textbf{Starter files:} \texttt{threads\_demo.c}, \texttt{banker.c}.\par
\textbf{Evidence to keep:} repeated counter results, critical-section explanation, Need matrix, and safe/unsafe trace.

\begin{labmode}
Threads and Banker's algorithm solve different problems. The first half asks how concurrent execution can safely share data; the second asks whether a resource-allocation state leaves a guaranteed route to completion.
\end{labmode}

\medskip\noindent\textbf{Learning outcomes.}
\begin{enumerate}
  \item explain why threads in one process can race on a shared variable;
  \item identify the critical section protected by a mutex;
  \item calculate Need and execute the Banker's safety test without confusing ``unsafe'' with ``already deadlocked.''
\end{enumerate}

\medskip\noindent\textbf{Part A --- Establish the correct threaded result.}
\begin{lstlisting}[style=osbash]
cd experiments/Lab06
gcc -Wall -Wextra threads_demo.c -o threads_demo -pthread
./threads_demo
\end{lstlisting}
Run the locked version three times. The final counter should equal the printed expected value. Find \texttt{pthread\_create}, \texttt{pthread\_join}, and the mutex lock/unlock pair.

\medskip\noindent\textbf{Part B --- Remove protection and observe a race.}
\begin{lstlisting}[style=osbash]
./threads_demo --no-lock
./threads_demo --no-lock
./threads_demo --no-lock
\end{lstlisting}
Record the actual counter each time. If one run happens to equal the expected value, do not conclude the program is safe; repeat it. Explain the lost-update idea using the read-modify-write sequence on \texttt{counter}.

\begin{keyidea}
A mutex protects a rule about shared state: only one thread at a time may execute the critical update. It does not make all thread execution sequential.
\end{keyidea}

\medskip\noindent\textbf{Part C --- Deadlock conditions as a design checklist.}
Consider two threads: T1 locks A then asks for B; T2 locks B then asks for A. Identify where mutual exclusion, hold-and-wait, no preemption, and circular wait appear. Then state one simple prevention rule based on lock ordering.

\medskip\noindent\textbf{Part D --- Compute Need before running Banker's algorithm.}
For each process calculate
\[
\texttt{Need} = \texttt{Maximum} - \texttt{Allocation}.
\]
Then compile and run \texttt{banker.c}. For every chosen process, record Work before, Need, and Work after its allocated resources are conceptually released.

\medskip\noindent\textbf{Part E --- Compare a safe and unsafe starting point.}
\begin{lstlisting}[style=osbash]
gcc -Wall -Wextra banker.c -o banker
./banker
./banker --unsafe
\end{lstlisting}
Explain the first point at which the unsafe run cannot make progress. Use the phrase ``no complete safe sequence can be guaranteed'' rather than claiming that an actual circular wait has already been observed.

\begin{verification}
Your record should include at least three locked and three unlocked counter results, a short critical-section explanation, the complete Need matrix, and a trace showing why the safe run progresses while the unsafe run stops.
\end{verification}

\begin{tryit}
Optional extension: change the initial Work vector by one resource unit and identify the exact Need $\leq$ Work comparison that becomes possible or impossible.
\end{tryit}
\end{labproject}
